% =====================================================================
\section{Especificación de requisitos final}
\label{sec:ers}

\subsection{Estructura del documento conforme a ISO/IEC/IEEE 29148:2018}

\begin{table}[H]
\centering\footnotesize
\caption{Correspondencia entre las secciones del ERS y ISO/IEC/IEEE 29148:2018}
\label{tab:estructura-ers}
\begin{tabular}{|C{1.2cm}|L{5.0cm}|L{4.2cm}|C{3.6cm}|}
\hline
\rowcolor{grisclaro}\textbf{§} & \textbf{Sección del ERS} & \textbf{Cláusula de 29148} & \textbf{Extensión}\\ \hline
1 & Introducción, propósito y alcance & 9.5.2 Introduction & 7 páginas\\ \hline
2 & Descripción general y partes interesadas & 9.5.3 Overall description & 12 páginas\\ \hline
3 & Requisitos específicos & 9.5.4 Specific requirements & 34 páginas\\ \hline
4 & Modelado UML e i* & 9.5.5 Verification (apoyo) & 11 páginas\\ \hline
5 & Priorización y trazabilidad & 6.5 Requirements traceability & 12 páginas\\ \hline
6 & Prototipo mínimo viable y limitaciones & 9.5.6 Supporting information & 6 páginas\\ \hline
7 & Conclusiones y trabajo pendiente & --- & 4 páginas\\ \hline
8 & Diagramas de flujo de datos & 9.5.5 (incorporada en la PE5) & 5 páginas\\ \hline
9 & Componentes de inteligencia artificial & 9.5.4 (incorporada en la PE5) & 8 páginas\\ \hline
\end{tabular}
\end{table}

\subsection{Catálogo final de requisitos}

\begin{table}[H]
\centering\footnotesize
\caption{Composición del catálogo de requisitos del ERS final}
\label{tab:catalogo}
\begin{tabular}{|L{5.6cm}|C{2.0cm}|C{2.4cm}|L{4.4cm}|}
\hline
\rowcolor{grisclaro}\textbf{Tipo} & \textbf{Cantidad} & \textbf{Rango} & \textbf{Observación}\\ \hline
Requisitos funcionales & 42 & \id{RF-01}--\id{RF-42} & 24 \emph{Must}, 16 \emph{Should}, 2 \emph{Could}\\ \hline
Requisitos no funcionales & 19 & \id{RNF-01}--\id{RNF-19} & Mapeados a las nueve características de ISO/IEC 25010:2023\\ \hline
Restricciones de diseño & 10 & \id{RD-01}--\id{RD-10} & Impuestas por el entorno o la contraparte\\ \hline
Requisitos legales & 8 & \id{RL-01}--\id{RL-08} & Derivados de la LOPDP\\ \hline
Requisitos de IA & 18 & \id{RF-IA-01}--\id{RNF-IA-12} & Incorporados en la Unidad V\\ \hline
\textbf{Total auditado} & \textbf{61} & --- & RF + RNF; las restricciones y los requisitos legales se auditan aparte\\ \hline
\end{tabular}
\end{table}

\subsection{Requisitos clave y por qué lo son}

De los cuarenta y dos requisitos funcionales, cinco concentran la propuesta de valor del sistema y
son los que conviene poder defender uno por uno.

\begin{table}[H]
\centering\footnotesize
\caption{Requisitos clave del SIMPA y su justificación en la evidencia}
\label{tab:req-clave}
\begin{tabular}{|C{1.4cm}|L{4.2cm}|L{4.4cm}|L{4.4cm}|}
\hline
\rowcolor{grisclaro}\textbf{ID} & \textbf{Requisito} & \textbf{Evidencia que lo sustenta} & \textbf{Por qué es clave}\\ \hline
\id{RF-07} & Diagnóstico de plaga o enfermedad por imagen & \id{EV-02}: la asesoría técnica visita cada quince días & Es la funcionalidad que reduce la ventana de detección de dos semanas a inmediata\\ \hline
\id{RF-22} & Estimación del porcentaje de fruta verde con alerta preventiva & \id{EV-04}: la extractora penaliza a partir del 3\,\% & Traslada la decisión de cortar desde el juicio visual a una medición anterior al despacho\\ \hline
\id{RF-28} & Registro de avance por unidad de labor & \id{EV-05}, \id{EV-06} & Es la entrada de la que dependen la liquidación, el indicador de desempeño y los reportes\\ \hline
\id{RF-35} & Registro delegado del avance & \id{EV-12}: el 11,3\,\% no tiene teléfono inteligente & Sin él, una parte del personal quedaría fuera del sistema por su equipamiento\\ \hline
\id{RF-37} & Cálculo de la remuneración semanal & \id{EV-13}: el cálculo manual consume una jornada semanal & Es el proceso de mayor frecuencia y mayor costo administrativo de la organización\\ \hline
\end{tabular}
\end{table}

\subsection{Catálogo completo de requisitos funcionales}

\begin{longtable}{|M{1.8cm}|L{9.6cm}|M{2.4cm}|}
\caption{Catálogo de requisitos funcionales del ERS final, agrupados por bloque}
\label{tab:rf-completo}\\
\hline
\rowcolor{grisclaro}\textbf{ID} & \textbf{Nombre del requisito} & \textbf{MoSCoW}\\ \hline
\endfirsthead
\hline
\rowcolor{grisclaro}\textbf{ID} & \textbf{Nombre del requisito} & \textbf{MoSCoW}\\ \hline
\endhead
\multicolumn{3}{|>{\columncolor{grissuave}}l|}{\textbf{Bloque I - Gestión de la estructura productiva}}\\ \hline
\id{RF-01} & Autenticación y control de acceso por rol & Must\\ \hline
\id{RF-02} & Gestión de plantaciones y lotes & Must\\ \hline
\id{RF-03} & Gestión de personal y equipos de trabajo & Must\\ \hline
\id{RF-10} & Gestión de variedades de palma & Must\\ \hline
\multicolumn{3}{|>{\columncolor{grissuave}}l|}{\textbf{Bloque II - Planificación y seguimiento de labores}}\\ \hline
\id{RF-26} & Planificación semanal de labores con presupuesto & Must\\ \hline
\id{RF-27} & Arrastre de labor incompleta a la semana siguiente & Should\\ \hline
\id{RF-34} & Registro de solicitud de insumos asociada al plan & Could\\ \hline
\id{RF-36} & Catálogo codificado de labores con tarifa por unidad & Must\\ \hline
\id{RF-37} & Liquidación semanal de presupuesto contra ejecución & Must\\ \hline
\id{RF-38} & Registro de labor no presupuestada & Should\\ \hline
\id{RF-39} & Aprobación de la liquidación semanal & Should\\ \hline
\id{RF-35} & Registro delegado para personal sin dispositivo propio & Must\\ \hline
\multicolumn{3}{|>{\columncolor{grissuave}}l|}{\textbf{Bloque III - Registro operativo de campo}}\\ \hline
\id{RF-04} & Registro de labores agrícolas & Must\\ \hline
\id{RF-28} & Registro de avance por unidad de labor & Must\\ \hline
\id{RF-29} & Consulta de avance acumulado y estimación de remuneración & Should\\ \hline
\id{RF-05} & Registro de monitoreo fitosanitario & Must\\ \hline
\id{RF-30} & Reporte de incidencia desde campo con evidencia fotográfica y georreferencia & Must\\ \hline
\id{RF-31} & Lista de verificación de pendientes por visita & Should\\ \hline
\id{RF-06} & Seguimiento de etapas del cultivo & Should\\ \hline
\id{RF-40} & Exportación de los datos personales de la persona titular & Must\\ \hline
\id{RF-41} & Rectificación de datos con conservación en bitácora & Must\\ \hline
\id{RF-42} & Supresión de datos personales con disociación del histórico productivo & Must\\ \hline
\id{RF-11} & Registro de análisis de suelo y foliar & Should\\ \hline
\id{RF-17} & Registro y consulta de datos climáticos & Should\\ \hline
\multicolumn{3}{|>{\columncolor{grissuave}}l|}{\textbf{Bloque IV - Polinización asistida}}\\ \hline
\id{RF-13} & Registro del proceso de polinización & Must\\ \hline
\id{RF-14} & Conteo georreferenciado de flores polinizadas & Must\\ \hline
\id{RF-15} & Registro y visualización de recorridos del personal & Should\\ \hline
\id{RF-16} & Evaluación de desempeño del personal & Should\\ \hline
\multicolumn{3}{|>{\columncolor{grissuave}}l|}{\textbf{Bloque V - Diagnóstico asistido por inteligencia artificial}}\\ \hline
\id{RF-07} & Detección de plagas y enfermedades por análisis de imagen & Must\\ \hline
\id{RF-08} & Diagnóstico de deficiencias nutricionales por análisis de imagen & Must\\ \hline
\id{RF-09} & Ficha técnica de plaga con recomendación de control & Should\\ \hline
\id{RF-21} & Clasificación de madurez del racimo por desprendimiento & Must\\ \hline
\id{RF-12} & Generación de alertas tempranas & Must\\ \hline
\multicolumn{3}{|>{\columncolor{grissuave}}l|}{\textbf{Bloque VI - Calidad de la fruta y relación con la extractora}}\\ \hline
\id{RF-22} & Estimación y alerta preventiva de porcentaje de fruta verde & Must\\ \hline
\id{RF-23} & Registro de la calificación de calidad emitida por la extractora & Should\\ \hline
\id{RF-24} & Trazabilidad de la calificación hasta la cuadrilla responsable & Should\\ \hline
\id{RF-25} & Control de la ventana entre corte y entrega & Should\\ \hline
\id{RF-33} & Registro de la guía de remisión del despacho & Could\\ \hline
\multicolumn{3}{|>{\columncolor{grissuave}}l|}{\textbf{Bloque VII - Reportes, estimación e histórico}}\\ \hline
\id{RF-18} & Estimación de producción y rendimiento & Must\\ \hline
\id{RF-32} & Compartición de la estimación de cosecha con la planta extractora & Should\\ \hline
\id{RF-19} & Generación de reportes & Must\\ \hline
\id{RF-20} & Historial de actividades, diagnósticos y eventos & Should\\ \hline
\end{longtable}

\noindent
La distribución resultante ---24 \emph{Must}, 16 \emph{Should} y 2 \emph{Could}--- concentra el
57\,\% del catálogo en la categoría de cumplimiento obligatorio. Es una proporción alta y conviene
justificarla en lugar de presentarla como dato: tres de los \emph{Must} (\id{RF-40} a \id{RF-42})
implementan obligaciones legales que no admiten priorización, y cuatro más (\id{RF-36} a \id{RF-39})
cubren el proceso de liquidación semanal, que la organización ejecuta sin excepción todas las
semanas. Un catálogo con menos \emph{Must} habría exigido declarar opcional alguna de esas dos cosas.

\subsection{Requisitos no funcionales por característica de calidad}

\begin{longtable}{|M{1.8cm}|L{3.6cm}|L{8.4cm}|}
\caption{Requisitos no funcionales del ERS final y su característica de calidad}
\label{tab:rnf-completo}\\
\hline
\rowcolor{grisclaro}\textbf{ID} & \textbf{Característica} & \textbf{Descripción abreviada}\\ \hline
\endfirsthead
\hline
\rowcolor{grisclaro}\textbf{ID} & \textbf{Característica} & \textbf{Descripción abreviada}\\ \hline
\endhead
\id{RNF-01} & Adecuación funcional & La clasificación de madurez del racimo alcanza una exactitud >= 85\% y una exhaustividad >= 80\% sobre la clase ``verde'' en el conjunto de prueb\\ \hline
\id{RNF-02} & Adecuación funcional & La detección de plagas alcanza una exactitud >= 80\% sobre las cinco afectaciones más frecuentes del cultivo.\\ \hline
\id{RNF-03} & Eficiencia de desempeño & Las operaciones de registro y consulta responden en <= 3 s para el percentil 95, con hasta 20 sesiones concurrentes. El valor se deriva de la pob\\ \hline
\id{RNF-04} & Eficiencia de desempeño & El diagnóstico por imagen entrega resultado en <= 10 s con ancho de banda >= 5 Mbps.\\ \hline
\id{RNF-05} & Eficiencia de desempeño & El registro GPS captura la ubicación con precisión <= 10 m y período de muestreo <= 30 s.\\ \hline
\id{RNF-06} & Compatibilidad & La exportación hacia la planta extractora se emite en formato CSV o JSON conforme al esquema documentado, sin campos de identificación personal.\\ \hline
\id{RNF-07} & Compatibilidad & La aplicación móvil opera en Android 9.0 o superior; el panel web, en las dos últimas versiones estables de Chrome, Edge y Firefox.\\ \hline
\id{RNF-08} & Interacción de usuario & Una persona sin experiencia previa completa el registro de labor y el análisis de imagen con una tasa de éxito >= 90\% tras una inducción <= 10 m\\ \hline
\id{RNF-09} & Interacción de usuario & La totalidad de la interfaz se presenta en español, con etiquetas correspondientes a la terminología del glosario y sin abreviaturas técnicas no \\ \hline
\id{RNF-10} & Fiabilidad & El servicio de respaldo presenta una disponibilidad mensual >= 99\%, equivalente a una indisponibilidad <= 7 h 12 min (720 h 1\%).\\ \hline
\id{RNF-11} & Fiabilidad & La aplicación opera sin conexión y sincroniza el 100\% de los registros pendientes al restablecerse la conectividad, sin pérdida ni duplicación.\\ \hline
\id{RNF-12} & Fiabilidad & Ante pérdida de señal GPS, el sistema conserva el registro parcial y lo marca como de precisión degradada, sin descartar la jornada.\\ \hline
\id{RNF-13} & Seguridad & Las credenciales se almacenan mediante función de derivación de clave con sal; la totalidad del tráfico emplea TLS 1.2 o superior.\\ \hline
\id{RNF-14} & Seguridad & Los datos de geolocalización se cifran en reposo y su acceso queda registrado en bitácora con identificación de la persona usuaria, fecha y motiv\\ \hline
\id{RNF-15} & Mantenibilidad & El sistema presenta arquitectura modular con cobertura de pruebas automatizadas >= 60\% en la capa de lógica de negocio.\\ \hline
\id{RNF-16} & Interacción de usuario & Explicabilidad - véase la especificación ampliada en la §3.3.1.\\ \hline
\id{RNF-17} & Flexibilidad (reemplazabilidad) & La migración del servicio de inferencia a un proveedor alternativo no requiere modificar el cliente móvil, al mediar una interfaz de servicio des\\ \hline
\id{RNF-18} & Flexibilidad (adaptabilidad) & Los umbrales de madurez, de porcentaje de fruta verde y de ventana corte-entrega son parametrizables por variedad y por lote sin recompilar la ap\\ \hline
\id{RNF-19} & Seguridad física (safety) & Toda recomendación de control fitosanitario emitida por el sistema (RF-09) incluye, de forma no omisible, el equipo de protección personal exigid\\ \hline
\end{longtable}

\subsection{Restricciones de diseño y requisitos legales}

\begin{table}[H]
\centering\footnotesize
\caption{Restricciones de diseño y su origen}
\label{tab:rd}
\begin{tabular}{|C{1.6cm}|L{6.8cm}|L{6.2cm}|}
\hline
\rowcolor{grisclaro}\textbf{ID} & \textbf{Restricción} & \textbf{Origen}\\ \hline
\id{RD-01} & Cliente de campo sobre Android 9.0 o superior & El 88,7\,\% del personal dispone de teléfono inteligente propio (\id{EV-12})\\ \hline
\id{RD-02} & Operación sin conexión con sincronización diferida & El 74,2\,\% no tiene señal permanente en el lote (\id{EV-12})\\ \hline
\id{RD-05} & La inferencia se ejecuta dentro de la frontera del sistema & Restricción de conectividad y de costo de servicio externo\\ \hline
\id{RD-06} & Interfaz adaptada a competencia digital baja & Brecha declarada entre administración y personal de campo (\id{EV-07})\\ \hline
\id{RD-10} & Registro delegado para quien no dispone de dispositivo & El 11,3\,\% del personal (\id{EV-12})\\ \hline
\end{tabular}
\end{table}

\noindent
Se presentan cinco de las diez restricciones, las que condicionan decisiones de arquitectura; el
catálogo completo reside en la Tabla 52 del ERS. Las cinco comparten una propiedad que conviene
destacar: ninguna proviene de una preferencia técnica del equipo. Todas se derivan de un dato del
contexto que la elicitación midió.

\begin{table}[H]
\centering\footnotesize
\caption{Requisitos legales derivados de la LOPDP y su implementación}
\label{tab:rl}
\begin{tabular}{|C{1.4cm}|C{1.6cm}|L{5.6cm}|L{5.6cm}|}
\hline
\rowcolor{grisclaro}\textbf{ID} & \textbf{Artículo} & \textbf{Obligación} & \textbf{Requisito que la implementa}\\ \hline
\id{RL-04} & Art. 12 & Derecho de acceso de la persona titular a sus datos & \id{RF-40}, exportación por la propia titular\\ \hline
\id{RL-05} & Art. 13 & Derecho de rectificación & \id{RF-41}, corrección con conservación en bitácora\\ \hline
\id{RL-06} & Art. 14 & Derecho de eliminación & \id{RF-42}, supresión con disociación del histórico\\ \hline
\end{tabular}
\end{table}

\noindent
Los tres requisitos de la columna derecha no existían en la Entrega 2A. El ERS enunciaba las
obligaciones en su mapeo legal sin que ninguna función del sistema las satisficiera, y la inspección
de la Unidad IV lo registró como el defecto \id{DEF-12}. Su reparación exigía especificar tres
requisitos funcionales nuevos, de modo que no admitía corrección editorial y se tramitó ante el
Change Control Board como \id{RFC-03}.

\subsection{Requisitos no funcionales y modelo de calidad}

Los diecinueve requisitos no funcionales se organizan según las nueve características de calidad de
producto de ISO/IEC 25010:2023 \cite{iso25010}. La cobertura de las nueve se declaró en la Entrega 2A
y resultó no sostenerse: la inspección de la Unidad IV detectó que se empleaba ``Portabilidad'' como
característica de primer nivel ---cuando en la edición 2023 sus subcaracterísticas dependen de
Flexibilidad--- y que ninguna característica de Seguridad física tenía requisito asociado. Ambos
hallazgos se tramitaron ante el Change Control Board como \id{RFC-02} y se cerraron con la
reclasificación de \id{RNF-17} y la especificación de \id{RNF-19}.

El caso es ilustrativo de una diferencia que conviene retener: declarar conformidad con una norma y
aplicar su estructura correctamente son dos cosas distintas, y solo la segunda es verificable.
